%==============================================================================
% FICHAMENTO CRÍTICO — Artificial Intelligence in Oral Cancer Diagnosis:
% A Systematic Review and Meta-Analysis
% Conforme NBR 6023 (referências) e NBR 10520 (citações)
%==============================================================================
\section{Fichamento: \citeonline{li2024oralcancer}}

% --- 1. IDENTIFICAÇÃO ---
\subsection{Identificação}
\begin{tabular}{ll}
\textbf{Título:} & Diagnostic accuracy of artificial intelligence assisted
clinical imaging in the detection of oral potentially malignant disorders
and oral cancer: a systematic review and meta-analysis \\
\textbf{Autor(es):} & Li, JingWen; Kot, Wai Ying; McGrath, Colman Patrick;
Chan, Bik Wan Amy; Ho, Joshua Wing Kei; Zheng, Li Wu \\
\textbf{Periódico:} & International Journal of Surgery \\
\textbf{Ano:} & 2024 \\
\textbf{Volume:} & 110 \\
\textbf{Número:} & 8 \\
\textbf{Páginas:} & 5034--5046 \\
\textbf{DOI:} & \url{https://doi.org/10.1097/JS9.0000000000001469} \\
\textbf{Qualis:} & A1 (Medicina) \\
\textbf{Tipo:} & Revisão sistemática com meta-análise \\
\end{tabular}

% --- 2. OBJETIVO ---
\subsection{Objetivo}
Examinar a aplicação de algoritmos de inteligência artificial na detecção de
desordens orais potencialmente malignas (OPMDs) e lesões cancerosas orais, e
avaliar a variação de acurácia entre diferentes ferramentas de imagem
(fotografia clínica, autofluorescência e tomografia de coerência óptica — OCT)
utilizadas nesses processos diagnósticos.

% --- 3. METODOLOGIA ---
\subsection{Metodologia}
\textbf{Desenho do estudo:} Revisão sistemática com meta-análise registrada
no Research Registry (UIN: reviewregistry1809), seguindo as diretrizes
PRISMA-DTA e MOOSE. \\
\textbf{Amostra:} 17 estudos publicados, totalizando 24.430 amostras
(7.020 pacientes), incluindo 12 estudos transversais, 3 retrospectivos e
2 prospectivos, com biópsia e histopatologia como padrão-ouro em 70,59\%
dos estudos. \\
\textbf{Métricas principais:} Odds ratio diagnóstica (DOR), sensibilidade,
especificidade, valor preditivo negativo (VPN), área sob a curva SROC,
análise de subgrupos por modalidade de imagem. \\
\textbf{Validação:} Meta-regressão para comparar modalidades de imagem;
análise de heterogeneidade (I²); gráfico GOSH para identificação de outliers
influentes.

% --- 4. RESULTADOS PRINCIPAIS ---
\subsection{Resultados Principais}
\begin{itemize}
  \item DOR global de 68,438 (IC 95\%: 39,484--118,623; I² = 86\%) e área
        sob a curva SROC de 0,938, indicando excelente desempenho diagnóstico.
  \item Sensibilidade de 89,9\% (IC 95\%: 86,6--92,5\%), especificidade de
        89,2\% (IC 95\%: 85,1--92,2\%) e VPN de 89,5\% (IC 95\%: 85,1--92,7\%).
  \item Após exclusão de outliers com GOSH plot: DOR = 49,30; AUC = 0,877;
        sensibilidade = 90,5\%; especificidade = 87,0\%.
  \item A fotografia clínica apresentou a maior acurácia diagnóstica
        (OR = 77,772; IC 95\%: 25,832--234,152) entre as modalidades de
        imagem analisadas.
\end{itemize}

% --- 5. LIMITAÇÕES ---
\subsection{Limitações}
\begin{itemize}
  \item Heterogeneidade substancial entre os estudos (I² = 86\% para DOR),
        parcialmente reduzida mas não eliminada após análise GOSH.
  \item Apenas 5,9\% dos estudos realizaram validação externa, comprometendo
        a generalização dos resultados para contextos clínicos diversos.
  \item A maioria dos estudos (15 de 17) foi conduzida na Ásia, com
        predomínio indiano, limitando a representatividade étnica e geográfica.
  \item Ausência de padronização nos protocolos de aquisição de imagem e
        nos limiares diagnósticos entre os estudos primários.
\end{itemize}

% --- 6. CONTRIBUIÇÃO PARA O LIVRO ---
\subsection{Contribuição para o Livro}
Este artigo fundamenta diretamente o Capítulo 11, que aborda revisões
sistemáticas e meta-análises sobre IA para diagnóstico de câncer oral. Os
dados agrupados de sensibilidade (89,9\%) e especificidade (89,2\%) fornecem
as estimativas mais abrangentes disponíveis sobre o desempenho de diferentes
modalidades de imagem assistidas por IA. A subanálise por tipo de imagem
(fotografia clínica, autofluorescência e OCT) é particularmente relevante
para a discussão sobre a escolha da modalidade de imagem no SDD do livro,
destacando que a fotografia clínica, por sua acessibilidade e baixo custo,
pode ser a ferramenta mais adequada para programas de rastreamento
populacional de câncer oral.

% --- 7. ANÁLISE CRÍTICA ---
\subsection{Análise Crítica}
\begin{itemize}
  \item \textbf{Validade interna:} Alta -- Revisão registrada, busca
        sistemática em 4 bases, seleção duplo-cega independente, análise
        de qualidade com QUADAS-2 e GOSH plot para outliers.
  \item \textbf{Validade externa:} Média -- Incluiu estudos multicêntricos
        asiáticos, mas com baixa representatividade de populações ocidentais
        e africanas.
  \item \textbf{Reprodutibilidade:} Alta -- Estratégia de busca e critérios
        de inclusão/exclusão detalhados e reproduzíveis a partir da
        publicação.
  \item \textbf{Relevância clínica:} Alta -- VPN de 89,5\% sugere que a IA
        pode reduzir significativamente o número de biópsias desnecessárias
        em programas de triagem.
  \item \textbf{Conflito de interesses:} Não -- Os autores declararam
        ausência de conflitos de interesse.
  \item \textbf{Viés identificado:} Viés de publicação -- heterogeneidade
        elevada (I² $>$ 80\%) e concentração geográfica sugerem possível
        viés de seleção. Apenas 5,9\% dos estudos reportaram validação
        externa, indicando viés de validação.
\end{itemize}

% --- 8. CITAÇÕES RELEVANTES ---
\subsection{Citações Relevantes}
\begin{citacao}
``The overall DOR for AI-based screening of OPMD and oral mucosal cancerous
lesions from normal mucosa was 68.438 (95\% CI = [39.484--118.623], I² = 86\%).
The area under the SROC curve was 0.938, indicating an excellent diagnostic
accuracy.'' (p. 5039).
\end{citacao}

\begin{citacao}
``Clinical photograph [OR = 77.772, 95\% CI (25.832--234.152)] demonstrated
the highest diagnostic accuracy, followed by OCT [OR = 56.997, 95\% CI
(24.540--132.383)] and autofluorescence tool [OR = 33.908, 95\% CI
(16.323--70.440)].'' (p. 5041).
\end{citacao}

% --- 9. MAPEAMENTO SDD ---
\subsection{Mapeamento SDD}
\textbf{Problema:} Necessidade de estimativa robusta do desempenho
diagnóstico de modelos de IA para detecção de OPMD e câncer oral a partir
de imagens clínicas, considerando diferentes modalidades de aquisição. \\
\textbf{Entrada:} 17 estudos primários com 24.430 amostras, incluindo
fotografias clínicas, imagens de autofluorescência e OCT, com padrão-ouro
histopatológico. \\
\textbf{Saída:} DOR, sensibilidade, especificidade, VPN e AUC agrupados
por modalidade de imagem; DOR ajustado pós-exclusão de outliers. \\
\textbf{Critérios:} Busca em $\geq$ 4 bases; avaliação de qualidade
metodológica (QUADAS-2); identificação e tratamento de heterogeneidade
com GOSH plot; meta-regressão por subgrupo. \\
\textbf{Confiança:} 0,82 -- Revisão robusta e abrangente, mas limitada pela alta
heterogeneidade e concentração geográfica dos estudos incluídos.
\end{tabular}

% --- 10. REFERÊNCIA ABNT ---
\subsection{Referência ABNT}
\begin{verbatim}
LI, JingWen et al. Diagnostic accuracy of artificial intelligence assisted
clinical imaging in the detection of oral potentially malignant disorders
and oral cancer: a systematic review and meta-analysis. International
Journal of Surgery, London, v. 110, n. 8, p. 5034-5046, 2024. DOI:
10.1097/JS9.0000000000001469.
\end{verbatim}
